\section{About This Book}\label{app:about}

This handbook is a study and prayer companion for keeping Advent in the 1962
Roman Rite. It is not an official liturgical book, a critical edition, or a
substitute for your parish calendar, your confessor, or your diocesan
ordinance. Wherever local feasts, transfers, or obligations differ from the
universal pattern described here, follow the calendar and law of your own
church.

\subsection{The Liturgical Books and Translations}

The liturgical order followed throughout is the universal Roman Rite of 1962:
the \textit{Missale Romanum} issued by decree of the Sacred Congregation of
Rites on 23 June 1962, governed by Pope John XXIII's motu proprio
\textit{Rubricarum instructum} (25 July 1960) and the \textit{Codex Rubricarum}
promulgated on 26 July 1960 (effective 1 January 1961), together with the
Sacred Congregation of Rites' instruction \textit{De musica sacra et sacra
liturgia} (3 September 1958).

For the Latin Collects of the Mass, the text is the 1962 Roman wording checked
against the 1962 typical edition. Because English hand missals occasionally
depart from the Latin they print beside them, each Collect is paired with a
verified public-domain translation that renders its Latin accurately:
F.~X.~Lasance and F.~A.~Walsh, \textit{The New Roman Missal in Latin and
English} (revised edition, 1945) for the First, Third, and Fourth Sundays,
Ember Wednesday's first Collect, Ember Friday, Ember Saturday, the Immaculate
Conception, and the Vigil of the Nativity; and Dom Gaspar Lefebvre, O.S.B.,
\textit{Daily Missal with Vespers} (1925) for the Second Sunday of Advent
(where Lasance prints a Son-addressed ending beside a Father-addressed Latin
prayer) and for Ember Wednesday's second Collect (\textit{Festina, quaesumus}).

Because a 1961--1962 typical edition of the \textit{Breviarium Romanum} is not
among the books consulted directly here, the Latin Office texts---the Isaiah
Matins lesson headings, the seven Great O Antiphons, \textit{Creator alme
siderum}, \textit{Alma Redemptoris Mater}, \textit{O magnum mysterium}, and
\textit{O admirabile commercium}---are taken from the Ratisbon winter Breviary
of 1888 (\textit{Breviarium Romanum, pars hiemalis}, Pustet), paired with the
English of John, Marquess of Bute, \textit{The Roman Breviary} (Winter volume,
1908), and read strictly under the 1960 structural rubrics (notably the
simplification of Advent Sunday Matins to one nocturn of nine psalms and three
lessons, RGBR nn.~162~a, 220, and the omission of impeded Scriptural lessons,
RGBR n.~213). The Christmas Eve \textit{Kalenda} proclamation is taken from the
1902 Propaganda Fide \textit{Martyrologium Romanum} and Bute (1908), noting
their minor two-year discrepancy from the flood (2957 in Latin, 2959 in Bute).
The four-verse \textit{Rorate caeli} prose comes from the 1903 Solesmes
\textit{Manual of Gregorian Chant} and Dom Prosper Gu\'eranger's \textit{The
Liturgical Year: Advent} (translated by Dom Laurence Shepherd, 1870), where it
appears as a French Advent chant and Benediction prose rather than a Roman
Missal or Breviary text. The \textit{Angelus} reproduces Ambrose St.~John's
1910 English \textit{Raccolta}.

English Scripture quotations follow the Challoner revision of the Douay-Rheims
Bible (using the 1899 American Douay-Rheims text for the Gospel canticles in
Appendix~\ref{app:prayers}), and Latin Scripture follows the Clementine
Vulgate. In the Latin \textit{Benedictus} (Luke~1:70) in
Appendix~\ref{app:prayers}, the electronic \texttt{latVUC} feed contains a
digitization error (\textit{locutum}); the vocal text printed here follows the
received liturgical reading \textit{sicut locutus est}, verified against
medieval and liturgical editions (Migne PL~142 and Blunt, 1866). Psalm numbers
follow the Vulgate numbering used in the 1962 Roman books, with modern Hebrew
numbering added in brackets where the two differ (for example,
Psalm~24~[25], \textit{Ad te levavi}). Where the original Greek of Luke~2:7 and
23:53 is discussed, the readings were checked in the SBL Greek New Testament and
the Robinson--Pierpont Byzantine textform.

\paragraph{Textual and historical apparatus notes.}
Several manuscript and edition details omitted from the narrative chapters for
readability are recorded here: (1) In St.~Gregory of Tours's \textit{Historia
Francorum} (X.31), the 1838 edition of Guadet and Taranne noted that their
\textit{codex Regius} omitted Bishop Perpetuus's St.~Martin's-Day-to-Christmas
fast, whereas Wilhelm Arndt's 1884 MGH edition showed that Paris lat.~10848
(codex~\textit{p}) preserves that clause while omitting the preceding
October-to-Martinmas clause. (2) The Second Council of Tours (567) monastic
fast canon is numbered c.~18 in Maassen's MGH edition and c.~17 in Hefele's
\textit{History of the Councils}; the First Council of M\^acon (c.~9) is dated
581 by Hefele, 582 by Gu\'eranger, and 583 by Maassen; and the Council of
Saragossa (c.~4) is dated 380 by the Spanish era~418 (379 in Mansi). (3) The
Old Gelasian Sacramentary rests on \textit{Vat.\ Reg.\ lat.\ 316} (edited by
H.~A.~Wilson, 1894); the Mozarabic manuscript vs.\ printed count is established
by Dom Marius F\'erotin (1912); the seven Roman Great O Antiphons appear in
order in the St.~Gall Hartker Antiphonary (\textit{Cod.\ Sang.\ 390}); and on
the Octave Day of Christmas, Bute renders \textit{Sicut pluvia in vellus} with
``dew'' rather than literal ``rain.''

The stations of the First Sunday, the Vigil of the Nativity, and the three
Masses of Christmas (St.~Mary Major at the Crib at Midnight, St.~Anastasia at
Dawn, and St.~Mary Major by Day) were checked directly in page images of the
1962 typical Missal and the 1862 Ratisbon Missal; the stations of the Second,
Third, and Fourth Sundays and the Ember Days rest on the 1862 Ratisbon and
1945 Lasance text layers.

\subsection{The Seven Calendar Shapes of Advent}

Under the 1960 \textit{Codex Rubricarum} (RG n.~20), the First Sunday of Advent
is the Sunday falling on 30 November or nearest to it---equivalently, the fourth
Sunday strictly before 25 December. Advent therefore opens on one of seven dates
from 27 November through 3 December, lasting from twenty-eight down to
twenty-two calendar days (liturgically, from First Vespers on Saturday evening
through None inclusive of 24 December, RG n.~71). The table below maps all
seven recurring shapes, verified across the century 2020--2120.

{\footnotesize
\begin{longtable}{>{\raggedright\arraybackslash}p{0.10\linewidth}>{\raggedright\arraybackslash}p{0.08\linewidth}>{\raggedright\arraybackslash}p{0.11\linewidth}>{\raggedright\arraybackslash}p{0.14\linewidth}>{\raggedright\arraybackslash}p{0.16\linewidth}>{\raggedright\arraybackslash}p{0.31\linewidth}}
\toprule
\textbf{Advent I} & \textbf{Days} & \textbf{Sun II / III} & \textbf{Ember W/F/S} & \textbf{Sun IV (O Ant.)} & \textbf{Key 1962 Overlaps \& Sample Years}\\
\midrule
\endhead
27 Nov & 28 & 4 / 11 Dec & 14 / 16 / 17 Dec & 18 Dec (\textit{O Adonai}) & Ember Sat on 17 Dec (\textit{O Sapientia}); 2022, 2033, 2039.\\
28 Nov & 27 & 5 / 12 Dec & 15 / 17 / 18 Dec & 19 Dec (\textit{O Radix}) & Ember Fri on 17 Dec, Sat on 18 Dec; 2021, 2027, 2032.\\
29 Nov & 26 & 6 / 13 Dec & 16 / 18 / 19 Dec & 20 Dec (\textit{O Clavis}) & St.~Lucy omitted on Sun III; 2020, 2026, 2037.\\
30 Nov & 25 & 7 / 14 Dec & 17 / 19 / 20 Dec & 21 Dec (\textit{O Oriens}) & St.~Andrew omitted on Sun I; St.~Thomas omitted on Sun IV; 2025, 2031, 2036.\\
1 Dec & 24 & 8 / 15 Dec & 18 / 20 / 21 Dec & 22 Dec (\textit{O Rex}) & Immaculate Conception wins over Sun II (priv.~comm.); St.~Thomas wins over Ember Sat (priv.~comm.); 2024, 2030, 2041.\\
2 Dec & 23 & 9 / 16 Dec & 19 / 21 / 22 Dec & 23 Dec (\textit{O Emmanuel}) & St.~Thomas wins over Ember Fri (priv.~comm.); 2029, 2035, 2040.\\
3 Dec & 22 & 10 / 17 Dec & 20 / 22 / 23 Dec & 24 Dec (Vigil) & Sun III on 17 Dec (\textit{O Sapientia}); Vigil replaces Sun IV with no commemoration (RG n.~30~a); 2023, 2028, 2034.\\
\bottomrule
\end{longtable}
}

In every shape, Week I ends by 9 December and Week II ends by 16 December, while
Ember Saturday and the Fourth Sunday always lie inside 17--24 December. The
evergreen course in Appendix~\ref{app:daily-course} therefore moves from the
weekly tables to the fixed 17--24 December table on 17 December without losing a
day in any year.

\subsection{Canons and Legal Authorities Consulted}

The legal statements in Section~\ref{sec:penance-law} distinguish three strata
by body of law, territory, and date (checked as of 25 November 2026):

\begin{itemize}
\item \textbf{Earlier historical discipline:} Second Council of Tours (567), c.~17 (18 in MGH), binding monks to daily December fasting; First Council of M\^acon (581/582/583), c.~9, establishing the Gallican Monday--Wednesday--Friday fast from St.~Martin's Day to Christmas; Council of Trent, Session XXIV, \textit{De reformatione matrimonii}, c.~10 (11 November 1563), forbidding solemn nuptials from Advent I through Epiphany; and the Third Plenary Council of Baltimore's \textit{Manual of Prayers} (New York, 1889, pp.~18--19), recording nineteenth-century United States practice.
\item \textbf{1962 universal Latin discipline:} \textit{Codex Iuris Canonici} (1917), c.~1108 \S\S\,1--3 (solemn nuptial blessing forbidden from Advent I through Christmas Day inclusive, unless permitted by the local Ordinary for just cause), c.~1252 \S\,1 (Friday abstinence), c.~1252 \S\,2 (fast and abstinence on the Advent Ember Days and the Vigil of Christmas, with no general Advent fast), and c.~1252 \S\,4 (cessation on Sundays and holy days of precept outside Lent without vigil anticipation); modified universally by the Sacred Congregation of the Council's decree of 3 December 1959 (AAS 51:918), permitting anticipation of the Christmas Vigil fast and abstinence from 24 December to 23 December. Because the 1917 Code's age-bound canon (c.~1254) and holy-day list (c.~1247) were not directly checked in this work's tracked legal records, those age limits are not stated.
\item \textbf{Current universal Latin law and United States particular law:} Pope St.~Paul VI's apostolic constitution \textit{Paenitemini} (17 February 1966, Norms II, IV, VI \S\,1~B), which removed statutory fast and abstinence from Ember Days and vigils; the \textit{Codex Iuris Canonici} (promulgated 25 January 1983, effective 27 November 1983), cc.~1246--1247 (Sunday Mass and ten universal holy days of obligation, including 8 December and 25 December, subject to episcopal-conference adaptation under c.~1246 \S\,2) and cc.~1249--1253 (statutory penitential days limited to Fridays and Lent, not Advent; abstinence from age 14; fasting from majority to the beginning of the sixtieth year); and for the United States, the National Conference of Catholic Bishops' \textit{Pastoral Statement on Penance and Abstinence} (18 November 1966, effective First Sunday of Advent, 27 November 1966, nos.~7, 17, 22--25) together with the USCCB Complementary Norm to cc.~1252--1253 (approved September 1983, promulgated 21 October 1983; fasting age amended 29 February 1984 to completion of the eighteenth year through the beginning of the sixtieth year). Because the 1991/1992 U.S. bishops' complementary norm on holy days of obligation (c.~1246 \S\,2) was not directly checked in tracked legal records, readers should verify annual holy-day obligation rules (especially when 8 December falls on a Saturday, Sunday, or Monday) with their own diocese or bishops' conference.
\end{itemize}

\subsection{Rights and Provenance}

All quoted Scripture, patristic translations, historical decrees, and received
prayers remain in the public domain or under their own rights as noted in
Appendix~\ref{app:references}; contemporary restricted texts of the Holy See
and the United States Conference of Catholic Bishops are summarized faithfully
with brief bounded phrases rather than reproduced in full. The
research scope, claim-level source audit, source records, and fresh-reader
review are maintained alongside the manuscript files in the repository.
